%%%%%%%%%%%%%%%%%%%%%%%%%%%%%%%%%%%%%%%%%%%%%%%%%%%%%%%%%%%%%%%%%%%%%%
%  第 3 个逻辑月 · 月度报告
%  主题：形态 · 物理
%
%  说明：组长对此月只给出「形态」「物理」两个关键词，
%        具体所指【待确认3】。本报告按「形态定型 + 物理基础可信度」
%        的理解组织，相关判断待组长确认。
%  编译：xelatex 本文件（跑两遍）
%%%%%%%%%%%%%%%%%%%%%%%%%%%%%%%%%%%%%%%%%%%%%%%%%%%%%%%%%%%%%%%%%%%%%%
\documentclass[stage2,fontset=fandol]{ydzjreport}

\stagetwosetup
  {第 3 个逻辑月 · 形态与物理}%
  {冯旭东、陈思文、张昊洋}%
  {复合微通道分区冷却}%
  {【待确认】逻辑月 3 对应的自然月}%
  {形态定型；单柱基准；物性复算；网格无关性}%
  {把多柱整体对比收敛到单柱基准，剥离单参数影响；
   并推进新物性 UDF 的挂接与代表工况复算，补齐物理基础。}%

\begin{document}
\makestagetwohead

\begin{notebox}

\begin{itemize}
  \item \textbf{翼型扰流柱形态参数}
  \item \textbf{物理推导辅助机器学习}
\end{itemize}

\end{notebox}

% ==================================================================
\section{本月研究任务}

\subsection{任务目标}

\begin{enumerate}
  \item 建立\textbf{单个扰流柱的基准模型}，对参数变量进行分组，逐一控制变量考察独立影响；
   \item 收尾\textbf{网格无关性}检验，明确后续建模的网格基准；
  \item 开展流场与温度场的\textbf{机理分析}，解释已观察到的结构差异。
\end{enumerate}

\subsection{完成情况}

\begin{table}[H]
\centering
\caption{本月任务完成情况}
\label{tab:tasks}
\begin{tabularx}{\textwidth}{@{}Y c Y@{}}
\toprule
\rowcolor{ydzjlight} 任务 & 状态 & 说明 \\
\midrule
网格无关性检验 & 阶段完成 & 完成七档对比；约 300 万局部接近，联合判据待复核 \\
单柱基准模型 & 部分完成 & 建立分组参数，采用单根独立仿真 \\
流场与温度场机理分析 & 部分完成 & 已定位高弯度翼型的尾缘流动死区 \\
\bottomrule
\end{tabularx}
\end{table}

% ==================================================================
\section{研究进展}

\subsection{已完成的工作}

\textbf{（一）水物性 UDF 的修正}

上月核查发现原水物性函数存在误差。本月完成\textbf{新物性 UDF 的编写与编译}：
采用 101325 Pa 下 REFPROP WATER.FLD v9.1 液态水数据拟合，
提供随温度变化的密度、黏度、导热系数与比热；
\textbf{比热与显焓采用同一多项式及其积分}，保证二者定义一致。
源码对流体温度 280 至 360 K 进行范围检查，并提供 300 K 与 330 K 的物性自检输出。

工作由「发现误差」推进到「完成新函数编译」，
后续还需完成\textbf{挂接、物性校验与代表工况复算}。

\begin{table}[H]
\centering
\caption{原水物性在 298.15 K 校验点的偏差}
\label{tab:udf}
\begin{tabularx}{\textwidth}{@{}Y c Y@{}}
\toprule
\rowcolor{ydzjlight} 物性 & 相对参考值偏差 & 影响 \\
\midrule
密度 & 偏低约 60.8\% & 直接影响流量与压降 \\
比热 & 偏低约 10.5\% & 影响温升与热量收支 \\
黏度 & 偏低约 26.7\% & 影响雷诺数与流动阻力 \\
\bottomrule
\end{tabularx}
\end{table}

把物性修正放在本节最前，是因为\textbf{物性偏差会同时改变流量、压降与换热量级}，
在其不可信的前提下讨论结构优劣没有意义。

\textbf{（二）网格无关性检验}

本节以 NACA2416 的 $N_{\mathrm{ref}}=3\,235\,081$ 单元（以下简称323万档）
为\textbf{唯一对比参考}，统一分析七档网格的温度、压降、换热与流阻指标，
并补充流动场及出口截面图。参考档是本阶段的比较基准，不作为精确解。

\textbf{1. 统一比较条件与误差定义。}

七档结果统一取 $t=0.2\,\mathrm{s}$，最终时间步长均为 $0.001\,\mathrm{s}$。
461万档早期保留了一段 $0.0001\,\mathrm{s}$ 的推进历史，且其边界层设置和网格尺度策略
与323万档存在差别，比较时保留这些条件差异。
平均温度沿用全域单元算术平均，压降取液体入口与出口面积加权静压之差。
对所有指标统一计算
\begin{equation}
e_{\phi,i}=\frac{|\phi_i-\phi_{323}|}{|\phi_{323}|}\times100\%,
\qquad \phi\in\{\overline T,T_{\max},\Delta p,Nu,h_c,f_D\}.
\label{eq:mesh_error}
\end{equation}
式~\eqref{eq:mesh_error} 中分母始终取323万档对应指标，
温度与压降三项差异均不超过 $1\%$ 作为本轮联合比较线。
参考档自身差异为零仅由定义决定，不能据此宣布网格无关性通过。

\begin{table}[H]\centering
\caption{NACA2416 在 $0.2\,\mathrm{s}$ 的七档网格结果}
\label{tab:mesh}
\begin{tabular}{rrrr}\toprule
单元数 / 万 & $\overline T$ / K & $T_{\max}$ / K & $\Delta p$ / Pa \\
\midrule
26.02 & 304.6684 & 344.6755 & 2505.204 \\
57.51 & 304.4113 & 345.2657 & 2389.274 \\
107.36 & 303.9449 & 345.1536 & 2440.766 \\
152.49 & 305.0933 & 346.0488 & 2389.938 \\
294.33 & 304.3736 & 345.9640 & 2355.838 \\
323.51 & 304.3009 & 345.9838 & 2348.320 \\
461.08 & 305.9860 & 346.6188 & 2448.714 \\
\bottomrule\end{tabular}\end{table}

表~\ref{tab:mesh} 保留七档原始值，图~\ref{fig:month3_mesh} 展示其变化；
曲线沿用已有分析图，收敛判断以离散数据点为准。
\begin{figure}[H]\centering
\includegraphics[width=\textwidth]{figs/month3_mesh_temperature_pressure.png}
\caption{七档网格的平均温度、最高温度与静压差原始值}
\label{fig:month3_mesh}\end{figure}

\textbf{2. 全部相对323万档的定量比较。}

\begin{table}[H]\centering
\caption{以323万档为参考的温度与压降相对差异（\%）}
\label{tab:mesh_ref323}
\begin{tabular}{rrrr}\toprule
单元数 / 万 & 平均温度差异 & 最高温度差异 & 压降差异 \\
\midrule
26.02 & 0.1208 & 0.3782 & 6.6807 \\
57.51 & 0.0363 & 0.2076 & 1.7440 \\
107.36 & 0.1170 & 0.2400 & 3.9367 \\
152.49 & 0.2604 & 0.0188 & 1.7722 \\
294.33 & 0.0239 & 0.0057 & 0.3201 \\
323.51 & 0.0000 & 0.0000 & 0.0000 \\
461.08 & 0.5537 & 0.1835 & 4.2751 \\
\bottomrule\end{tabular}\end{table}

按表~\ref{tab:mesh_ref323}，294万档相对323万档的平均温度、最高温度和压降差异分别为
$0.0239\%$、$0.0057\%$ 和 $0.3201\%$，三项均低于 $1\%$。
两档绝对差分别为 $0.0726\,\mathrm{K}$、$0.0198\,\mathrm{K}$ 和 $7.518\,\mathrm{Pa}$，
表明294--323万附近存在局部接近区间。152万档压降差异为 $1.7722\%$，
不能将约150万单元认定为本组已验证的网格基准。
461万档相对323万档的平均温度、最高温度和压降差异分别为
$0.5537\%$、$0.1835\%$ 和 $4.2751\%$，压降仍超过比较线。
因此，统一以323万为参考后，局部两档接近与全部网格达到同一阈值仍需区分。

图~\ref{fig:month3_errors323} 将温度、压降和换热流阻指标的参考口径统一到323万档。
流体侧换热采用 $D_h=2.8070\,\mathrm{mm}$、$L=3.5\,\mathrm{mm}$；
物性取入口、出口质量流量加权温度平均值对应的数值。
323万档 $Nu=175.818$、$h_c\approx38210.0\,\mathrm{W/(m^2\cdot K)}$，
294万档相对其差异均约为 $0.11\%$。
461万档 $Nu=173.114$、$h_c\approx37621.5\,\mathrm{W/(m^2\cdot K)}$，
相对323万档差异均约为 $1.54\%$；$f_D$ 差异约为 $4.28\%$。
$Re$ 与 $Pr$ 接近主要反映入口及物性条件接近，不能替代换热与压降核查。

\begin{figure}[H]\centering
\includegraphics[width=0.98\textwidth]{figs/month3_mesh_errors_ref323.png}
\caption{各指标统一以323万档为参考的相对差异；虚线为 $1\%$ 比较线}
\label{fig:month3_errors323}\end{figure}

\textbf{3. 流动场与出口回流的图像核查。}

图~\ref{fig:month3_flowfield} 以323万档为参照，对照461万档在
$z=1.50\,\mathrm{mm}$ 平面的主流向速度、平面投影流线、温度与静压。
流线是三维速度在该平面的投影，黑色边界为 $U_x=0$，不能直接等同于完整三维流线。
两档均呈现翼型间主流加速、尾缘低速区和沿流向的压力降低；
局部逆向速度区的形状与范围存在差别。
原始图中两档使用相同变量色标，便于按323万档观察局部差异，
但二维图像本身不能证明空间离散误差已收敛。

\begin{figure}[H]\centering
\includegraphics[width=\textwidth]{figs/month3_mesh_flowfield.png}
\caption{以323万档（上排）为参照的中高平面流动场对比；下排为461万档，$t=0.2\,\mathrm{s}$}
\label{fig:month3_flowfield}\end{figure}

图~\ref{fig:month3_outlet} 进一步给出出口法向速度与温度分布。
323万档的回流面积比例为 $5.1720\%$，回流质量流量占入口流量的 $0.4754\%$；
461万档对应为 $5.3886\%$ 与 $0.2777\%$。
相对323万档，461万档回流面积比例增加约0.2167个百分点，
回流质量流量比例反而减少约0.1977个百分点，不能据此称其回流整体更严重。
两档回流面质量加权温度均为 $305\,\mathrm{K}$。
出口局部回流与内部尾迹应分别识别，不能只凭投影流线将二者混为一类。

\begin{figure}[H]\centering
\includegraphics[width=\textwidth]{figs/month3_mesh_outlet.png}
\caption{以323万档（上排）为参照的出口速度与温度对比；下排为461万档}
\label{fig:month3_outlet}\end{figure}

以相同统计区域的体积加权温度辅助核查时，323万档与461万档分别为
$305.7793\,\mathrm{K}$ 和 $305.8673\,\mathrm{K}$，
相对323万档仅约 $0.0288\%$。
这说明单元算术平均值受网格统计权重影响；七档主表仍保留原始统一口径，
两档体积加权数值单列说明，不能与其他档算术平均值混用。
流场与出口对比能定位差异位置，但未进行出口延长或回流温度敏感性重算，
不能把 $4.2751\%$ 的压降差异直接归因于某一种出口因素。

\textbf{4. 以323万为细档的 GCI 与阶段结论。}

GCI 沿用以323万为细档、从较粗网格选取对比档的已有结果，
采用三维等效尺度 $\ell=(V/N)^{1/3}$、安全系数 $F_s=1.25$。
294万与323万的细化比仅约1.032，尺度间隔较近。
107万、152万、323万组合的压降形式 GCI 为 $0.7275\%$，表观阶约5.57；
26万、107万、323万组合的 $Nu$ 与 $h_c$ 形式 GCI 约为 $7.13\%$ 与 $7.24\%$。
GCI 是三档网格外推指标，不等于式~\eqref{eq:mesh_error} 中的两档相对差异，
不能通过挑选最小的形式 GCI 认定全部指标通过。

本轮统一以323万档作为\textbf{后续诊断与成本折中的参考网格}，
294万档相对该参考的温度、压降和换热指标均较接近；
但461万档的压降与换热差异仍需解释，现有结果不支持全部网格均已通过联合检验。
其他构型资料另行保留，直线与混合结构的网格质量合格只表明准备条件满足，
gao2.1 较早批次也不与本节七档序列混用。
后续应统一初场、时间推进历史和网格策略，并以323万档为参照，
继续核查出口、近壁及尾迹区的敏感性。原始数据、统一误差表与流场出处均保留在源码目录。

\textcolor{red}{王佳蔚和大树哥哥补充图片说明}

\textbf{（三）机理分析：高弯度翼型 PEC 回落的原因}

\textbf{1. 先明确评价口径。}

局部对流换热满足牛顿冷却定律：

\begin{equation}
q'' = h\,(T_w - T_f),
\label{eq:newton}
\end{equation}

其中 $q''$ 为壁面热流密度，$T_w$ 为壁面温度，$T_f$ 为流体温度。
努塞尔数与水力直径定义为

\begin{equation}
Nu = \frac{h\,D_h}{k_f},\qquad D_h = \frac{4A_c}{P_w},
\label{eq:nu}
\end{equation}

$k_f$ 为流体导热系数，$A_c$ 为流通截面积，$P_w$ 为湿周。
达西摩擦系数由进出口压差反算：

\begin{equation}
f = \frac{2\,\Delta p\,D_h}{\rho u_m^2 L},
\label{eq:friction}
\end{equation}

$u_m$ 为截面平均流速，$L$ 为流道长度。
综合性能取\textbf{等泵功基准}下的常用定义：

\begin{equation}
PEC = \frac{Nu/Nu_0}{\left(f/f_0\right)^{1/3}},
\label{eq:pec}
\end{equation}

下标 0 表示基准构型（NACA0015）。
式~\eqref{eq:pec} 的含义是：分子衡量换热增益，分母衡量水力代价，
二者都相对于同一基准，因此 PEC 可以横向比较不同构型。

\textbf{2. 再追机理链条。}

翼型弯度增大后，前缘迎风侧流速提高，而尾缘背风侧的扩压段变长、
逆压梯度增强：

\begin{equation}
\frac{\mathrm{d}p}{\mathrm{d}x} > 0 .
\label{eq:adverse}
\end{equation}

当逆压梯度超过边界层动量的储备，边界层发生\textbf{分离}。
分离后壁面附近形成回流区，壁面法向速度梯度趋于零：

\begin{equation}
\left.\frac{\partial u}{\partial y}\right|_{w} \rightarrow 0 ,
\label{eq:shear}
\end{equation}

即壁面剪切率显著下降。

对流换热系数与\textbf{壁面温度梯度}直接相关：

\begin{equation}
h = -\,\frac{k_f}{T_w - T_f}\left.\frac{\partial T}{\partial y}\right|_{w},
\label{eq:hgrad}
\end{equation}

这可由式~\eqref{eq:newton} 与傅里叶导热定律在壁面处的匹配直接得到。
速度梯度下降使热边界层增厚，$\left|\partial T/\partial y\right|_w$ 随之减小，
因此 $h$ 在死区附近\textbf{显著下降}。

\textbf{3. 两个后果，方向一致。}

\begin{enumerate}
  \item \textbf{有效换热面积减小。}
  死区内 $h$ 很低，等效于这部分壁面基本不参与换热；
  总对流热阻
  $R_{\mathrm{conv}} = 1/(hA_{\mathrm{eff}})$ 上升。
  \item \textbf{流阻同时上升。}
  回流区占据流通截面，主流被挤压、局部速度提高，$f$ 上升。
\end{enumerate}

代入式~\eqref{eq:pec} 可以看到两者对 PEC 的作用方向\textbf{完全一致}：
分子 $Nu/Nu_0$ 下降、分母 $(f/f_0)^{1/3}$ 上升，PEC 必然回落。
这就解释了「弯度越大换热越好」的预期为何不成立——
\textbf{弯度带来的增益被尾缘分离抵消，还额外付出了流阻}。

\textbf{4. 与厚度效应的对照。}

厚度增加的作用机理不同：它主要提高迎风侧扰动强度与换热面积，
\textbf{并未引起同等规模的分离}，因此 $Nu$ 上升而 $f$ 增幅有限，
PEC 得以提高。这正是 NACA0024、NACA4424 的 PEC 高于基准、
而 NACA6424、NACA6624 回落的原因。

\textbf{5. 本轮结论的边界。}

上述分析是\textbf{定性}的。要给出定量判断，还缺三个量：
分离点位置、回流区长度、壁面局部 $Nu$ 分布。
这三项都需要从单柱模型或流场后处理中提取，列入下月工作。

\textbf{（四）参数分组设置与变量独立检验}
\textcolor{red}{王佳蔚陈思文补充分组变量和后续基础图片和一定数据}

% ==================================================================
\section{问题与对策}

\subsection{遇到的主要问题}

\begin{enumerate}
  \item \textbf{物性修正后尚未复算。}
        新 UDF 已完成编写与编译，但还没有挂接到算例上，
        因此无法回答「修正物性后结构排序会不会变」这个问题。
  \item \textbf{单柱基准模型尚未建立。}
        整体流道内多柱干扰无法解耦，厚度、弯度、攻角三个参数的独立影响仍不能分离。
  \item \textbf{机理分析停留在定性。}
        尾缘流动死区已定位，但分离点位置、回流区长度与局部 $Nu$ 分布都还没有定量结果。
  \item \textbf{本月关键词的工作范围待确认。}
        「形态」「物理」两个词具体指什么，需要组长明确后才能定稿。
\end{enumerate}

\subsection{已尝试的方案}

\begin{enumerate}
  \item 针对问题 1，已完成新 UDF 的编写、编译与温度范围检查；
        挂接与复算需要在下月继续。
  \item 针对问题 2，已从流场分析入手定位现象原因（尾缘流动死区），
        属于定性解释；定量剥离仍需单柱模型。
  \item 针对问题 3，本报告按"物性 + 机理"两者都做的理解组织，
        待组长确认后可调整范围。
\end{enumerate}

\subsection{下一步对策}

\begin{enumerate}
  \item 挂接新 UDF，完成物性自检与代表工况复算，确认排序是否变化；
  \item 建立单柱基准模型，采用与整体流道一致的网格基准与边界条件；
  \item 在单柱模型上做单参数扫掠，得到厚度、弯度、攻角的独立影响曲线。
\end{enumerate}

% ==================================================================
\section{阶段性结论}

\subsection{可以确定的}

\begin{enumerate}
  \item 已完成 NACA2416 七档网格结果整理，294 万与 323 万档局部接近；
        统一以323万为参考时，294万档温度与压降差异低于 $1\%$；
        461万档压降差异为 $4.2751\%$，参考网格的适用范围仍待复核。
  \item 原水物性函数存在量级误差（密度偏低约 60.8\%），已识别并完成新 UDF 编译。
  \item 高弯度翼型 PEC 回落的原因已定位为\textbf{尾缘流动死区}导致有效换热面积减小；
        其机理已从牛顿冷却定律与壁面温度梯度关系展开到分离判据
        （式~\eqref{eq:newton}--\eqref{eq:hgrad}）。
  \item \textbf{物理基础必须先于结构比较完成}——
        在物性不可信的情况下比较结构，等于在错误的基准上排序。
\end{enumerate}

\subsection{尚不能确定的}

\begin{enumerate}
  \item \textbf{物性修正对结论的实际影响幅度。}
        缺少的证据是：挂接新 UDF 后代表工况的复算结果。
  \item \textbf{单柱条件下三个翼型参数的独立影响权重。}
        缺少的证据是：单柱基准模型上的单参数扫掠数据。
  \item \textbf{尾缘分离的定量特征。}
        缺少的证据是：分离点位置、回流区长度与壁面局部 $Nu$ 分布——
        有了这三项，式~\eqref{eq:hgrad} 的定性链条才能变成定量预测。
  \item \textbf{「形态」与「物理」的最终工作范围。}
        取决于组长对本月关键词的明确。
\end{enumerate}

% ==================================================================
\section{下月计划}

\begin{table}[H]
\centering
\caption{下月工作计划}
\label{tab:plan}
\begin{tabularx}{\textwidth}{@{}Y Y c@{}}
\toprule
\rowcolor{ydzjlight} 事项 & 预期产出 & 时间 \\
\midrule
挂接新 UDF，完成物性自检与代表工况复算 & 复算结果与排序对比 & 第 1 周 \\
建立单柱基准模型 & 单柱模型与基准数据 & 第 1--2 周 \\
单参数扫掠（厚度、弯度、攻角） & 参数影响特性曲线 & 第 2--4 周 \\
小块连接段设计（为下一逻辑月铺垫） & 连接段形态方案 & 第 3--4 周 \\
\bottomrule
\end{tabularx}
\end{table}

% ==================================================================
\begin{thebibliography}{9}
\bibitem{udf} 项目组. water\_properties\_refprop\_fit.c：液态水物性拟合与 Fluent UDF 源码. 项目计算资料, 2026.
\bibitem{iapws} IAPWS. SR6-08(2011)：0.1 MPa 下液态水物性补充发布. 2011.
\bibitem{monitor} Ansys. Fluent 用户指南：收敛监测. 2025 R1.
\end{thebibliography}

\end{document}
